% =====================================================================
%  CARNET D'ENTRAÎNEMENT — Fiche 12 (Chimie organique)
%  THÈME 4 — Masse molaire d'une molécule organique
%  NIVEAU MOYEN — CORRIGÉ
% =====================================================================
\documentclass[11pt,a4paper]{article}
\usepackage[utf8]{inputenc}
\usepackage[T1]{fontenc}
\usepackage{lmodern}
\usepackage[french]{babel}
\usepackage{amsmath,amssymb,mathtools}
\providecommand{\degree}{^\circ}
\usepackage{icomma}
\usepackage{siunitx}
\sisetup{output-decimal-marker={,},per-mode=symbol}
\usepackage[version=4]{mhchem}
\usepackage{chemfig}
\setchemfig{atom sep=2em,bond length=0.9em,cram width=2.5pt}
\usepackage{xcolor}
\usepackage{array}
\usepackage{booktabs}
\usepackage{tabularx}
\usepackage[most]{tcolorbox}
\usepackage{enumitem}
\usepackage{tikz}
\usetikzlibrary{arrows.meta,positioning,calc}
\usepackage{fancyhdr}
\usepackage{titlesec}
\usepackage[hidelinks]{hyperref}
\usepackage[a4paper,margin=2.1cm,headheight=26pt,headsep=9pt]{geometry}

\definecolor{primary}{RGB}{150,98,162}
\definecolor{primarydark}{RGB}{92,52,110}
\definecolor{excol}{RGB}{0,135,90}
\definecolor{attncol}{RGB}{200,40,55}
\definecolor{methcol}{RGB}{90,90,100}

\newcommand{\runtitle}{Thème 4 — Moyen — Corrigé}
\pagestyle{fancy}\fancyhf{}
\fancyhead[L]{\small\textcolor{primarydark}{\textbf{Fiche 12} — Chimie organique}}
\fancyhead[R]{\small\textcolor{primarydark}{\runtitle}}
\fancyfoot[C]{\small\thepage}
\renewcommand{\headrulewidth}{0.8pt}
\renewcommand{\headrule}{\hbox to\headwidth{\color{primary}\leaders\hrule height \headrulewidth\hfill}}
\setlength{\parindent}{0pt}\setlength{\parskip}{3pt}

\tcbset{boxbase/.style={enhanced,breakable,boxrule=0.9pt,arc=2.5pt,
   left=3mm,right=3mm,top=2mm,bottom=2mm,fonttitle=\bfseries,coltitle=white,
   toptitle=0.5mm,bottomtitle=0.5mm}}
\newtcolorbox[auto counter]{corr}[1]{boxbase,colback=excol!5,colframe=excol,
   title={\textsc{Corrigé}~\thetcbcounter\ \textmd{--- \itshape #1}}}
\newcommand{\rep}[1]{\textcolor{primarydark}{\textbf{#1}}}

\newcommand{\banner}[2]{%
\begin{tcolorbox}[enhanced,colback=excol,colframe=excol,arc=5pt,boxrule=0pt,
   left=5mm,right=5mm,top=2.5mm,bottom=2.5mm,coltext=white]
{\large\bfseries #1}\\[1pt]{\small #2}
\end{tcolorbox}\vspace{2pt}}

\begin{document}

\banner{Thème 4 — Masse molaire d'une molécule organique}{Niveau \textbf{Moyen} — Corrigé détaillé \quad$\vert$\quad Fiche 12, Chimie organique}

% ---------------------------------------------------------------------
\begin{corr}{masse molaire du paracétamol}
\[ M(\ce{C8H9NO2}) = \underbrace{8\times12}_{96} + \underbrace{9\times1}_{9} + \underbrace{1\times14}_{14} + \underbrace{2\times16}_{32} = 96+9+14+32 = \rep{\qty{151}{\g\per\mol}}. \]
\end{corr}

% ---------------------------------------------------------------------
\begin{corr}{nombre de molécules (aspirine)}
On convertit $\qty{360}{\mg}=\qty{0,360}{\g}$, puis :
\[ n=\frac{m}{M}=\frac{\qty{0,360}{\g}}{\qty{180}{\g\per\mol}}=\qty{2,0e-3}{\mol}, \qquad
N=n\,N_A=\qty{2,0e-3}{}\times\qty{6,02e23}{}\approx\rep{\qty{1,2e21}{}\ \text{molécules}}. \]
\end{corr}

% ---------------------------------------------------------------------
\begin{corr}{concentration d'une solution (paracétamol)}
$\qty{755}{\mg}=\qty{0,755}{\g}$ et $\qty{200}{\mL}=\qty{0,200}{\L}$.
\[ n=\frac{m}{M}=\frac{\qty{0,755}{\g}}{\qty{151}{\g\per\mol}}=\qty{5,0e-3}{\mol}, \qquad
C=\frac{n}{V}=\frac{\qty{5,0e-3}{\mol}}{\qty{0,200}{\L}}=\rep{\qty{2,5e-2}{\mol\per\L}}. \]
\end{corr}

% ---------------------------------------------------------------------
\begin{corr}{quelle masse peser ? (mannitol)}
\textbf{a)} $M(\ce{C6H14O6}) = 6\times12 + 14\times1 + 6\times16 = 72+14+96 = \rep{\qty{182}{\g\per\mol}}$.\\
\textbf{b)} On passe des molécules à la quantité de matière, puis à la masse :
\[ n=\frac{N}{N_A}=\frac{\qty{1,20e21}{}}{\qty{6,02e23}{}}\approx\qty{2,0e-3}{\mol}, \quad
m=n\,M=\qty{2,0e-3}{}\times\qty{182}{}\approx\qty{0,364}{\g}=\rep{\qty{364}{\mg}}. \]
\end{corr}

% ---------------------------------------------------------------------
\begin{corr}{le piège de l'unité}
\textbf{a)} $M_r = 6\times12+14\times1+2\times14 = 72+14+28 = \rep{114}$ (sans unité).\\
\textbf{b)} \rep{Non}, c'est faux : «\,$\si{\g\per\mol}$\,» est l'unité de la masse \emph{molaire} $M$, pas de la masse moléculaire \emph{relative} $M_r$ qui est un \textbf{nombre sans dimension}. Réponse correcte : $M_r = \rep{114}$ (et, si l'on demandait la masse molaire, $M = \qty{114}{\g\per\mol}$).
\end{corr}

\end{document}
